% VI21-FULL-SOLUTION-REFINEMENT: P18
\subsection*{VI.21.P18: Why local stalk maps are required}
\begin{problem}\label{prob:vi21-p18}
Explain precisely why the local-homomorphism condition on stalks is essential in the definition of a scheme morphism.
\end{problem}

\ifdefined\FullProblemDossiers
\paragraph{Why this problem matters.}
The local condition is the mechanism that makes field-valued points and residue fields functorial.

\paragraph{Useful ideas before starting.}
Translate point statements through local rings and residue fields; when the target is affine, use the VI/19 universal property and then verify the pointwise interpretation.

\paragraph{Guided hints.}
\begin{enumerate}
\item Identify the relevant canonical ring or stalk homomorphism.
\item Track maximal ideals or prime contractions to identify the image point.
\item Use locality to pass to residue fields when needed.
\item Check independence, uniqueness, or functoriality explicitly rather than only on points.
\end{enumerate}

\begin{solution}
A morphism of ringed spaces supplies maps
\[
f_x^\#:\OO_{Y,f(x)}\to\OO_{X,x},
\]
but without locality these maps need not be compatible with the distinguished maximal ideals of
the two local rings.  The scheme condition requires
\[
(f_x^\#)^{-1}(\mathfrak m_x)=\mathfrak m_{f(x)}.
\]
This has two crucial consequences.  First, the stalk map descends canonically to a field map
$\kappa(f(x))\to\kappa(x)$.  Second, it forces the sheaf map to correspond to the prescribed
continuous image point: elements vanishing at $f(x)$ pull back to elements vanishing at $x$, and
units at $f(x)$ pull back to units at $x$.

If locality were dropped, a ring map between stalks could send the maximal ideal incorrectly, so
its induced prime/closed-point behavior could disagree with the continuous map and no residue-field
map need exist.  The locally ringed condition is therefore exactly the compatibility tying
geometry of points to pullback of functions.
\end{solution}

\paragraph{What happened mathematically.}
The local condition is the mechanism that makes field-valued points and residue fields functorial.

\paragraph{Extensions.}
Relate this to the uniqueness formula for maps into affine schemes from VI/19.

\paragraph{Provenance.}
\textbf{Canonical VI/21 derivative/application; not an additional legacy migration.}
\else
\paragraph{Provenance.} \textbf{Canonical VI/21 derivative/application; not an additional legacy migration.}
\fi
